\centerline{\bf \Huge Налоговые вычеты 2}

\begin{flushright}
        \scriptsize{Опыт - это то, что получаешь, не получив того, что хотел.\\Айзен Соуске}
\end{flushright}

\textbf{\large Несложные задачи}

\problem{1}
Докажите, что если простое число $p$ является делителем числа $x^2+ 33x+22$, где $x$ - целое, то оно также является делителем числа $y^2+70y+224$ для некоторого целого $y$.

\problem{2}
Докажите, что если простое число $p$ является делителем числа $x^2-x+3$, где $x$ - целое, то оно также является делителем числа $y^2-y+25$ для некоторого целого $y$.

\problem{3}
Вычислите символы Лежандра:

1) $\left(\frac{111}{541}\right)$;

2) $\left(\frac{529}{601}\right)$;

3) $\left(\frac{2108}{2003}\right)$;

4) $\left(\frac{19525}{1847}\right)$.

\problem{4}
Опишите все простые числа $p$, для которых разрешимы сравнения:

(а) $x^2+1 \equiv 0 \pmod p;$

(б) $x^2-2 \equiv 0 \pmod p;$

\problem{5}
С помощью пункта предыдущей задачи докажите, что существует бесконечно много простых:

(а) вида $p \equiv 1\pmod 4$

(б) вида $p \equiv -1\pmod 8$


\textbf{\large Сложные задачи}

\problem{6}
Пусть $p>2$ - простое число, $a, b \in \mathbb{Z}$, причём $p \nmid a$. Докажите, что

$$
\sum_{x=0}^{p-1}\left(\frac{a x+b}{p}\right)=0 .
$$


\newpage

\hspace{3mm}
    
\problem{7}
Докажите, что уравнение $(4 m-1)(4 n-1)=4 x^2+1$ не имеет решений в натуральных числах.


\textbf{\large Невероятно сложные задачи}

\problem{8}
Let $m, n \geq 3$ be positive odd integers. Prove that $2^m-1$ doesn't divide $3^n-1$.


\problem{9}
Find all positive integers $n$ such that for any integer $k$ there exists an integer $a$ for which $a^3+a-k$ is divisible by $n$.

%https://artofproblemsolving.com/community/c6t385f6h582821_x3x_is_a_surjection_mod_n